\chapter{Controle: Esperar e Repetir}

\begin{objetivos}
\begin{itemize}
\item Fazer o programa esperar e controlar o tempo;
\item Usar \textbf{repita}, \textbf{sempre} e \textbf{repita até};
\item Parar o programa na hora certa.
\end{itemize}
\end{objetivos}

\section{Esperando}

Os blocos laranja \bloco[scControl]{Controle} comandam o \textbf{ritmo} do programa. O mais simples é o \bloco[scControl]{espere \entrada{1} segundos}, que faz o programa dar uma pausa antes de seguir.

\begin{dica}
Use pausas pequenas, como \entrada{0.1} ou \entrada{0.5}, para criar animações suaves.
\end{dica}

\section{Três jeitos de repetir}

\begin{center}
\renewcommand{\arraystretch}{1.8}
\begin{tabular}{@{}ll@{}}
\bloco[scControl]{repita \entrada{10} vezes} & repete um número exato de vezes \\
\bloco[scControl]{sempre} & repete para sempre, até clicar em parar \\
\bloco[scControl]{repita até que \booleano{...}} & repete até uma condição ficar verdadeira \\
\end{tabular}
\end{center}

\begin{pensaassim}[Pense assim: três tipos de repetição na vida real]
\begin{itemize}
\item \textbf{repita 10 vezes}: ``dê 10 pulinhos''. Você sabe \emph{quantas} vezes vai fazer.
\item \textbf{sempre}: o seu \emph{coração}, que bate sem parar enquanto você está vivo.
\item \textbf{repita até que}: ``continue caminhando \emph{até} chegar na porta''. Você não sabe quantos passos serão, só sabe \emph{quando parar}.
\end{itemize}
\end{pensaassim}

\begin{codigo}[Os três laços]
\begin{pilha}
\bk[scEvents]{quando \bandeira\ for clicado}
\sbc[scControl]{repita \entrada{4} vezes}
\bk[scMotion]{mova \entrada{10} passos}
\sbcend
\end{pilha}
\begin{pilha}
\bk[scEvents]{quando \bandeira\ for clicado}
\sbc[scControl]{sempre}
\bk[scMotion]{mova \entrada{10} passos}
\sbcend
\end{pilha}
\end{codigo}

\begin{vocabulario}[Vocabulário: laço infinito]
Um laço com \bloco[scControl]{sempre} se chama \textbf{laço infinito}: ele nunca termina sozinho. É perfeito para jogos, em que algo precisa ser vigiado o tempo todo.
\end{vocabulario}

\begin{cuidado}[Cuidado: laço sem fim]
Tudo o que fica \textbf{dentro} do laço se repete; o que fica \textbf{depois} dele só acontece quando o laço termina. Como o \bloco[scControl]{sempre} nunca termina, nenhum bloco colocado embaixo dele será executado!
\end{cuidado}

\begin{missao}[Missão 7: Gato guarda-noturno]
Faça o gatinho andar de um lado para o outro do palco, sem parar.
\end{missao}

\begin{codigo}[Patrulha]
\begin{pilha}
\bk[scEvents]{quando \bandeira\ for clicado}
\sbc[scControl]{sempre}
\bk[scMotion]{mova \entrada{10} passos}
\bk[scMotion]{se tocar na borda, volte}
\sbcend
\end{pilha}
\end{codigo}

\begin{cuidado}
Dentro de um laço, o gatinho pode ficar de cabeça para baixo ao voltar! Para evitar, clique no gatinho e em \textbf{Direção}, escolha a opção ``só esquerda/direita'', ou use o bloco \bloco[scMotion]{defina o estilo de rotação \menu{esquerda-direita}}.
\end{cuidado}

\section{Repetindo até uma condição}

O bloco \bloco[scControl]{repita até que} repete até que a pergunta lá dentro seja respondida com \textbf{verdadeiro}. Ele é como dizer ``continue andando \emph{até} chegar na parede''.

\begin{codigo}[Andar até a borda do palco]
\begin{pilha}
\bk[scEvents]{quando \bandeira\ for clicado}
\sbc[scControl]{repita até que \booleano{tocando em \menu{borda}?}}
\bk[scMotion]{mova \entrada{5} passos}
\sbcend
\end{pilha}
\end{codigo}

O sensor de borda detecta o limite do palco. Para detectar uma parede desenhada, veja o labirinto do capítulo 19.

\section{Parando}

\bloco[scControl]{pare \menu{tudo}} encerra o programa na hora. Use quando o jogo acabar.

\section{Hora de praticar}

\begin{exercicio}[Contagem regressiva]
Faça o gatinho dizer ``3'', ``2'' e ``1'' (cada um por 1 segundo) e depois ``Já!''.
\end{exercicio}

\begin{solucao}
\begin{pilha}
\bk[scEvents]{quando \bandeira\ for clicado}
\bk[scLooks]{diga \entrada{3} por \entrada{1} segundos}
\bk[scLooks]{diga \entrada{2} por \entrada{1} segundos}
\bk[scLooks]{diga \entrada{1} por \entrada{1} segundos}
\bk[scLooks]{diga \entrada{Já!} por \entrada{2} segundos}
\end{pilha}
\end{solucao}

\begin{exercicio}[Gato girador]
Faça o gatinho girar sem parar, 5 graus para a direita a cada 0.05 segundo.
\end{exercicio}

\begin{solucao}
\begin{pilha}
\bk[scEvents]{quando \bandeira\ for clicado}
\sbc[scControl]{sempre}
\bk[scMotion]{gire \entrada{5} graus}
\bk[scControl]{espere \entrada{0.05} segundos}
\sbcend
\end{pilha}
\end{solucao}

\begin{exercicio}[Cinco batidas]
Adicione a extensão \textbf{Música}. Faça o gatinho tocar o tambor 5 vezes. Depois de cada batida, espere meio segundo. Essa pausa começa quando o bloco de tambor termina.
\end{exercicio}

\begin{solucao}
\begin{pilha}
\bk[scEvents]{quando \bandeira\ for clicado}
\sbc[scControl]{repita \entrada{5} vezes}
\bk[scSound]{toque o tambor \menu{(1) Caixa} por \entrada{0.25} tempos}
\bk[scControl]{espere \entrada{0.5} segundos}
\sbcend
\end{pilha}
\end{solucao}

\begin{desafio}[Corrida até a borda do palco]
Faça o gatinho correr (mudando de fantasia a cada passo) até encostar na borda do palco e então dizer ``Cheguei!''.
\end{desafio}

\begin{solucao}
\begin{pilha}
\bk[scEvents]{quando \bandeira\ for clicado}
\sbc[scControl]{repita até que \booleano{tocando em \menu{borda}?}}
\bk[scMotion]{mova \entrada{5} passos}
\bk[scLooks]{próxima fantasia}
\bk[scControl]{espere \entrada{0.1} segundos}
\sbcend
\bk[scLooks]{diga \entrada{Cheguei!} por \entrada{2} segundos}
\end{pilha}
\end{solucao}
